\section{第4章 不定积分}
\subsection{不定积分的概念与性质}
\begin{exercise}{424}{基础·原函数关系}
设 $f$ 在 $(1,+\infty)$ 上的一个原函数为 $F(x)=\ln(x+\sqrt{x^2-1})$。求 $f(x)$。
\end{exercise}
\begin{exercise}{425}{基础·原函数与常数}
已知 $\ln x$ 是 $f$ 在 $(0,+\infty)$ 上的一个原函数。求 $f(x)$，并判断 $\ln(2x)$ 是否也是它的原函数。
\end{exercise}
\begin{exercise}{426}{常规·两个原函数之差}
在 $(0,1)$ 上求 $F(x)=\arcsin(2x-1)$、$G(x)=2\arcsin\sqrt x$ 的导数，说明它们为何相差常数，并求这个常数。
\end{exercise}
\begin{exercise}{427}{常规·分段原函数}
设 $f(x)=3x^2$（$x>0$），$f(x)=2x$（$x\le0$）。验证 $F(x)=x^3$（$x>0$）、$F(x)=x^2$（$x\le0$）是 $f$ 在实轴上的原函数，单独检查 $x=0$。
\end{exercise}
\begin{exercise}{428}{基础·连续函数的原函数}
说明 $f(x)=|x|$ 在实轴上存在原函数的依据，并验证 $F(x)=x|x|/2$ 是它的一个原函数。
\end{exercise}
\begin{exercise}{429}{常规·分段原函数拼接}
求 $f(x)=x$（$x<0$）、$f(x)=2x$（$x\ge0$）在实轴上满足 $F(0)=0$ 的原函数。
\end{exercise}
\begin{exercise}{430}{基础·原函数与导函数}
已知 $\sin x$ 是 $f(x)$ 的一个原函数。求 $f(x)$ 及 $\int f(x)\,\dd x$。
\end{exercise}
\begin{exercise}{431}{常规·斜率与初值}
一条曲线经过 $(e,-1)$，在每个 $x>0$ 处的切线斜率为 $\ln x/x$。求曲线方程。
\end{exercise}
\begin{exercise}{432}{基础·积分与微分}
设 $f$ 可导。分别求 $\int f'(x)\,\dd x$ 和 $\dd[\int f'(x)\,\dd x]$。
\end{exercise}
\begin{exercise}{433}{基础·积分与求导}
分别计算 $\int3x^2\,\dd x$、$\dd(\int3x^2\,\dd x)$ 及 $(\int3x^2\,\dd x)'$。
\end{exercise}
\begin{exercise}{434}{常规·原函数之差}
设 $F'=G'$ 在开区间 $I$ 上成立。用中值定理说明 $F-G$ 为常数；若 $F(x_0)=G(x_0)$，求这个常数。
\end{exercise}
\begin{exercise}{435}{基础·微分的积分}
求 $\int\dd(x^2+\sin x)$，并对结果求导核验。
\end{exercise}
\begin{exercise}{436}{常规·换元恢复导函数}
已知 $f'(\sin^2t)=\cos2t$（$0<t<\pi/2$），且 $f(1/2)=0$。先求 $f'(x)$（$0<x<1$），再求 $f(x)$。
\end{exercise}
\begin{exercise}{437}{基础·拆项积分}
求 $\int(x-1/x-3/\sqrt{1-x^2})\,\dd x$，在 $(-1,0)$ 或 $(0,1)$ 上分别考虑。
\end{exercise}
\begin{exercise}{438}{基础·三角恒等变形}
求 $\int\cos2x/(\sin^2x\cos^2x)\,\dd x$，在 $\sin x\cos x\ne0$ 的区间上考虑。提示：将分子写成 $\cos^2x-\sin^2x$。
\end{exercise}
\begin{exercise}{439}{基础·三角基本积分}
求 $\int(1-2\cos x)/\sin^2x\,\dd x$，在不含 $k\pi$ 的开区间上考虑。
\end{exercise}
\begin{exercise}{440}{基础·指数化简积分}
求 $\int[(5e)^{-x}+1-e^x]\,\dd x$。
\end{exercise}
\begin{exercise}{441}{基础·降幂积分}
求 $\int[6\sqrt x+5e^x+\cos^2(x/2)]\,\dd x$（$x>0$）。
\end{exercise}
\begin{exercise}{442}{常规·多项式除法积分}
求 $\int[x^4/(1+x^2)-2/\sqrt{1-x^2}]\,\dd x$（$|x|<1$）。提示：先对有理项作除法。
\end{exercise}
\begin{exercise}{443}{基础·对数原函数}
在 $(0,+\infty)$ 上求 $\int\dd x/x$，并求满足 $F(1)=2$ 的原函数。
\end{exercise}
\begin{exercise}{444}{常规·分段原函数小综合}
设 $f(x)=x$（$x<0$）、$f(x)=2x$（$x\ge0$）。求它在实轴上满足 $F(0)=1$ 的原函数，并检查 $F'(0)$。
\end{exercise}
\subsection{换元积分法}
\begin{exercise}{445}{基础·凑微分}
求 $\int\dd x/\sqrt{2x+3}$（$x>-3/2$）及 $\int(x+1)/(x^2+2x+5)\,\dd x$。
\end{exercise}
\begin{exercise}{446}{基础·根式换元}
求 $\int e^{2\sqrt x}/\sqrt x\,\dd x$ 与 $\int\sin\sqrt x/\sqrt x\,\dd x$（$x>0$）。提示：令 $t=\sqrt x$。
\end{exercise}
\begin{exercise}{447}{基础·复合函数换元}
求 $\int x^2e^{x^3}\,\dd x$ 和 $\int x/(1+x^4)\,\dd x$。
\end{exercise}
\begin{exercise}{448}{基础·三角换元}
求 $\int\sin^2x\cos^2x\,\dd x$ 与 $\int\sin^3x\cos x\,\dd x$。
\end{exercise}
\begin{exercise}{449}{基础·平方和与平方差积分}
求 $\int\dd x/(4+x^2)$ 和 $\int\dd x/(4-x^2)$，后一式在不含 $\pm2$ 的开区间上考虑。
\end{exercise}
\begin{exercise}{450}{基础·指数复合换元}
求 $\int e^{\tan x}\sec^2x\,\dd x$ 与 $\int(1-x^{-2})e^{x+1/x}\,\dd x$，分别在被积函数连续的开区间上考虑。
\end{exercise}
\begin{exercise}{451}{基础·对数换元}
求 $\int\dd x/[x\sqrt{1-\ln^2x}]$（$e^{-1}<x<e$）及 $\int(1+\ln x)/(1+x^2\ln^2x)\,\dd x$（$x>0$）。
\end{exercise}
\begin{exercise}{452}{基础·反三角函数换元}
求 $\int\arcsin x/\sqrt{1-x^2}\,\dd x$（$|x|<1$）和 $\int\arcsin\sqrt x/\sqrt{x(1-x)}\,\dd x$（$0<x<1$）。
\end{exercise}
\begin{exercise}{453}{常规·配方换元}
求 $\int\dd x/\sqrt{x-x^2}$（$0<x<1$）。提示：配方 $x-x^2=1/4-(x-1/2)^2$。
\end{exercise}
\begin{exercise}{454}{基础·三角内函数}
求 $\int\cos x/(9+\sin^2x)\,\dd x$，并求满足 $F(0)=0$ 的原函数。
\end{exercise}
\begin{exercise}{455}{基础·根式拆项积分}
求 $\int(1-x)/\sqrt{9-4x^2}\,\dd x$（$|x|<3/2$）。
\end{exercise}
\begin{exercise}{456}{常规·立方根换元}
求 $\int(\sin x+\cos x)/\sqrt[3]{\sin x-\cos x}\,\dd x$，立方根按实数定义，在 $\sin x-\cos x\ne0$ 的开区间上考虑。
\end{exercise}
\begin{exercise}{457}{常规·根式三角复合}
求 $\int x\tan(\sqrt{1-x^2})/\sqrt{1-x^2}\,\dd x$（$|x|<1$）。提示：令 $u=\sqrt{1-x^2}$。
\end{exercise}
\begin{exercise}{458}{常规·对数根式换元}
求 $\int[\ln(x+\sqrt{1+x^2})+1]/\sqrt{1+x^2}\,\dd x$。提示：令 $u=\ln(x+\sqrt{1+x^2})$。
\end{exercise}
\begin{exercise}{459}{常规·指数根式换元}
求 $\int e^{\sqrt{1+\sin x}}\cos x/\sqrt{1+\sin x}\,\dd x$，在 $1+\sin x>0$ 的开区间上考虑。提示：令 $u=\sqrt{1+\sin x}$。
\end{exercise}
\begin{exercise}{460}{基础·拆项与反正切换元}
求 $\int[x+(\arctan x)^2]/(1+x^2)\,\dd x$。
\end{exercise}
\begin{exercise}{461}{常规·三角幂有理化}
求 $\int\sin x\cos^3x/(1+\cos^2x)\,\dd x$。提示：令 $u=\cos x$，再作多项式除法。
\end{exercise}
\begin{exercise}{462}{常规·正切换元}
在 $(-\pi/2,\pi/2)$ 上求 $\int\dd x/(\sin^2x+2\cos^2x)$。提示：令 $t=\tan x$。
\end{exercise}
\begin{exercise}{463}{基础·指数换元}
求 $\int\dd x/(1+e^x)$，并求满足 $F(0)=0$ 的原函数。
\end{exercise}
\begin{exercise}{464}{常规·对数差整体换元}
求 $\int[\ln(x+1)-\ln x]/[x(x+1)]\,\dd x$（$x>0$）。提示：令 $u=\ln(x+1)-\ln x$。
\end{exercise}
\begin{exercise}{465}{基础·对数换元与拆分}
求 $\int\ln^2x/[x(1+\ln^2x)]\,\dd x$（$x>0$）。提示：令 $t=\ln x$。
\end{exercise}
\begin{exercise}{466}{常规·根式有理化}
求 $\int\dd x/(x-\sqrt{x^2-1})$（$x>1$）。提示：用共轭式将被积函数化为 $x+\sqrt{x^2-1}$。
\end{exercise}
\begin{exercise}{467}{基础·指数内函数凑微分}
求 $\int e^{\sin2x}\cos2x\,\dd x$。
\end{exercise}
\begin{exercise}{468}{常规·对数正切换元}
在 $0<x<\pi/2$ 上求 $\int\ln(\tan x)/\sin2x\,\dd x$。提示：计算 $\ln\tan x$ 的导数。
\end{exercise}
\begin{exercise}{469}{常规·商函数导数识别}
求 $\int(x\cos x-\sin x)/x^2\,\dd x$，在两个半轴上分别考虑。提示：计算 $(\sin x/x)'$。
\end{exercise}
\begin{exercise}{470}{常规·三角指数整体换元}
求 $\int e^{\sin x}\cos x/(1+e^{\sin x})\,\dd x$。提示：令 $u=e^{\sin x}$。
\end{exercise}
\begin{exercise}{471}{基础·半角积分}
求 $\int\sin x/(1+\cos x)\,\dd x$，在 $1+\cos x\ne0$ 的开区间上考虑。
\end{exercise}
\begin{exercise}{472}{常规·已知原函数换元}
已知 $\int f(u)\,\dd u=\ln(\sin u)+C$（$0<u<\pi$）。求 $\int xf(1-x^2)\,\dd x$，在 $|x|<1$ 上考虑。
\end{exercise}
\begin{exercise}{473}{常规·抽象函数商的积分}
设 $f\in C^1(I)$ 且 $0\notin I$。利用 $[f(x)/x]'$ 求 $\int[xf'(x)-f(x)]/x^2\,\dd x$。
\end{exercise}
\begin{exercise}{474}{常规·根式换元综合}
求 $\int x^3\sqrt{4-x^2}\,\dd x$（$|x|<2$）。提示：令 $u=4-x^2$，将余下的 $x^2$ 写成 $4-u$。
\end{exercise}
\begin{exercise}{475}{常规·平方和根式积分}
求 $\int x^2/(1+x^2)^{3/2}\,\dd x$。提示：拆成 $1/\sqrt{1+x^2}-1/(1+x^2)^{3/2}$。
\end{exercise}
\begin{exercise}{476}{常规·平方差根式积分}
求 $\int x/(x^2-4)^{3/2}\,\dd x$（$x>2$）。提示：令 $u=x^2-4$。
\end{exercise}
\begin{exercise}{477}{常规·立方根有理化}
求 $\int\dd x/[1+\sqrt[3]{x+1}]$（$x>-1$）。提示：令 $t=\sqrt[3]{x+1}$，将 $3t^2/(1+t)$ 作多项式除法。
\end{exercise}
\begin{exercise}{478}{常规·不同根指数换元}
求 $\int\dd x/(\sqrt x+\sqrt[3]{x^2})$（$x>0$）。提示：令 $t=x^{1/6}$，将积分化为 $\int6t^2/(1+t)\,\dd t$。
\end{exercise}
\begin{exercise}{479}{常规·根式分式换元}
求 $\int\dd x/[(x+1)\sqrt x]$（$x>0$）。提示：令 $t=\sqrt x$。
\end{exercise}
\begin{exercise}{480}{基础·指数有理积分}
求 $\int\dd x/(e^x+e^{2x})$。提示：令 $t=e^x$。
\end{exercise}
\begin{exercise}{481}{常规·指数根式换元}
求 $\int\dd x/\sqrt{e^x+1}$。提示：令 $t=\sqrt{e^x+1}$，则 $\dd x=2t\,\dd t/(t^2-1)$。
\end{exercise}
\begin{exercise}{482}{常规·幂换元与拆分}
求 $\int\dd x/[x(x^2+1)]$，分别在两个半轴上考虑。提示：拆成 $1/x-x/(x^2+1)$。
\end{exercise}
\begin{exercise}{483}{基础·倒数换元}
求 $\int\sin(1/x)/x^2\,\dd x$（$x\ne0$）。
\end{exercise}
\begin{exercise}{484}{常规·混合内函数换元}
求 $\int(1+x)/(x^2+x\ln x)\,\dd x$，在 $x>1$ 上考虑。提示：令 $u=x+\ln x$。
\end{exercise}
\begin{exercise}{485}{常规·配方与三角代换}
求 $\int\sqrt{2x-x^2}\,\dd x$（$0<x<2$）。提示：先配方为 $1-(x-1)^2$。
\end{exercise}
\begin{exercise}{486}{常规·指数根式复合}
求 $\int e^x/\sqrt{1+e^{2x}}\,\dd x$。提示：令 $t=e^x$。
\end{exercise}
\begin{exercise}{487}{常规·三角代换}
求 $\int x^2/\sqrt{1-x^2}\,\dd x$（$|x|<1$）。提示：令 $x=\sin t$，$|t|<\pi/2$，再降幂。
\end{exercise}
\begin{exercise}{488}{基础·换元积分变式}
求 $\int x/(1+3x^2)\,\dd x$。
\end{exercise}
\begin{exercise}{489}{常规·换元与初值综合}
求 $\int x/\sqrt{1+x^2}\,\dd x$，并求满足 $F(0)=2$ 的原函数。
\end{exercise}
\subsection{分部积分法}
\begin{exercise}{490}{基础·多项式乘指数或三角函数}
求 $\int xe^{2x}\,\dd x$ 和 $\int x^2\cos x\,\dd x$。
\end{exercise}
\begin{exercise}{491}{基础·反三角与对数分部积分}
求 $\int\ln x\,\dd x$（$x>0$）和 $\int\arctan x\,\dd x$。
\end{exercise}
\begin{exercise}{492}{常规·循环积分}
用两次分部积分求 $\int e^{-x}\cos x\,\dd x$。
\end{exercise}
\begin{exercise}{493}{基础·移位对数积分}
求 $\int x\ln(x+1)\,\dd x$（$x>-1$）。
\end{exercise}
\begin{exercise}{494}{基础·反正弦分部积分}
求 $\int\arcsin(x+1)\,\dd x$（$-2<x<0$）。
\end{exercise}
\begin{exercise}{495}{常规·凑微分与分部配合}
求 $\int x\cot x/\sin^2x\,\dd x$，在 $\sin x\ne0$ 的开区间上考虑。提示：$\cot x/\sin^2x=-(\cot^2x)'/2$。
\end{exercise}
\begin{exercise}{496}{基础·三角函数分部积分}
求 $\int x/\cos^2x\,\dd x$，在 $\cos x\ne0$ 的开区间上考虑。
\end{exercise}
\begin{exercise}{497}{常规·对数三角分部积分}
在 $0<x<\pi$ 上求 $\int\ln(\sin x)/\sin^2x\,\dd x$。提示：取 $u=\ln\sin x$、$\dd v=\csc^2x\,\dd x$。
\end{exercise}
\begin{exercise}{498}{常规·根式分部积分}
用分部积分求 $\int x^3/\sqrt{1+x^2}\,\dd x$。提示：取 $u=x^2$、$\dd v=x\,\dd x/\sqrt{1+x^2}$。
\end{exercise}
\begin{exercise}{499}{常规·指数三角循环积分}
用两次分部积分求 $\int e^{2x}\sin3x\,\dd x$。
\end{exercise}
\begin{exercise}{500}{常规·平方和根式积分}
求 $\int\sqrt{x^2+1}\,\dd x$。提示：分部后使用 $x^2=(x^2+1)-1$，将同类积分移项。
\end{exercise}
\begin{exercise}{501}{常规·平方差根式积分}
求 $\int\sqrt{4-x^2}\,\dd x$（$|x|<2$）。可令 $x=2\sin t$。
\end{exercise}
\begin{exercise}{502}{常规·反正弦乘多项式}
求 $\int x\arcsin x\,\dd x$（$|x|<1$）。提示：先取 $u=\arcsin x$，再对余下的根式积分作三角代换。
\end{exercise}
\begin{exercise}{503}{常规·反正切乘负幂}
求 $\int\arctan x/x^2\,\dd x$，在两个半轴上分别考虑。可使用 $1/[x(1+x^2)]=1/x-x/(1+x^2)$。
\end{exercise}
\begin{exercise}{504}{基础·对数乘有理因子}
求 $\int\ln x/x^2\,\dd x$（$x>0$）。
\end{exercise}
\begin{exercise}{505}{常规·指数换元与分部积分}
求 $\int xe^{\sqrt x}/\sqrt x\,\dd x$（$x>0$）。提示：令 $t=\sqrt x$，积分化为 $2\int t^2e^t\,\dd t$。
\end{exercise}
\begin{exercise}{506}{常规·对数与根式分部积分}
求 $\int\ln x/\sqrt x\,\dd x$（$x>0$）。提示：取 $u=\ln x$、$\dd v=x^{-1/2}\,\dd x$。
\end{exercise}
\begin{exercise}{507}{常规·反正切根式积分}
求 $\int\arctan\sqrt x\,\dd x$（$x>0$）。提示：分部积分后令 $t=\sqrt x$。
\end{exercise}
\begin{exercise}{508}{常规·指数有理导数识别}
求 $\int e^x(x-1)/x^2\,\dd x$（$x\ne0$）。提示：计算 $(e^x/x)'$。
\end{exercise}
\begin{exercise}{509}{常规·三角商的逆向求导}
求 $\int[(1+x)\sin x+\cos x]/(1+x)^2\,\dd x$（$x\ne-1$）。提示：计算 $[\cos x/(1+x)]'$。
\end{exercise}
\begin{exercise}{510}{常规·恢复导函数与分部积分}
已知 $f'(e^t)=\cos t$ 对所有实数 $t$ 成立且 $f(1)=0$。求 $f(x)$（$x>0$）。提示：先恢复 $f'(x)$，再令 $u=\ln x$。
\end{exercise}
\begin{exercise}{511}{常规·正割三次积分}
在 $(-\pi/2,\pi/2)$ 上用分部积分求 $\int\sec^3x\,\dd x$。提示：取 $u=\sec x$、$\dd v=\sec^2x\,\dd x$，再使用 $\tan^2x=\sec^2x-1$。
\end{exercise}
\begin{exercise}{512}{基础·分部积分变式}
求 $\int x\sin x\,\dd x$。
\end{exercise}
\begin{exercise}{513}{常规·对数平方积分}
求 $\int(\ln x)^2\,\dd x$（$x>0$），并确定满足 $F(1)=0$ 的原函数。
\end{exercise}
\subsection{有理函数的积分}
\begin{exercise}{514}{基础·假分式积分}
求 $\int x^3/(x^2+1)\,\dd x$。提示：先作多项式除法。
\end{exercise}
\begin{exercise}{515}{基础·两个一次因子}
求 $\int(3x-2)/(x^2+3x-4)\,\dd x$，在不含 $-4,1$ 的开区间上考虑。
\end{exercise}
\begin{exercise}{516}{常规·三个一次因子}
求 $\int\dd x/[(x+1)(x-1)(x+2)]$，在分母非零的开区间上考虑。
\end{exercise}
\begin{exercise}{517}{基础·简单部分分式}
求 $\int\dd x/[x(x-1)]$，在不含 $0,1$ 的开区间上考虑。
\end{exercise}
\begin{exercise}{518}{常规·重一次因子}
求 $\int(x^2+1)/[(x+1)^2(x-1)]\,\dd x$。提示：设分解为 $A/(x-1)+B/(x+1)+C/(x+1)^2$。
\end{exercise}
\begin{exercise}{519}{常规·一次与二次因子}
求 $\int\dd x/[(x^2+1)(1+2x)]$。提示：分解为 $A/(1+2x)+(Bx+C)/(x^2+1)$。
\end{exercise}
\begin{exercise}{520}{基础·配方与分子拆分}
求 $\int(3x+1)/(x^2+2x+5)\,\dd x$。
\end{exercise}
\begin{exercise}{521}{常规·重二次分母}
求 $\int\dd x/(x^2+1)^2$。提示：令 $x=\tan t$。
\end{exercise}
\begin{exercise}{522}{常规·四次分母分解}
求 $\int(2x^3+2x^2+5x+5)/(x^4+5x^2+4)\,\dd x$。提示：分母为 $(x^2+1)(x^2+4)$，可先把分子拆为奇、偶两部分。
\end{exercise}
\begin{exercise}{523}{常规·平移后重二次积分}
求 $\int(x+1)/(x^2+2x+2)^2\,\dd x$。提示：令 $u=x^2+2x+2$。
\end{exercise}
\begin{exercise}{524}{基础·高次式换元}
求 $\int x^2/(1+x^6)\,\dd x$。提示：令 $u=x^3$。
\end{exercise}
\begin{exercise}{525}{常规·重零点分解}
求 $\int\dd x/[x^2(1-x^2)]$，在不含 $-1,0,1$ 的开区间上考虑。提示：先拆为 $1/x^2+1/(1-x^2)$。
\end{exercise}
\begin{exercise}{526}{基础·重因子换元}
求 $\int x^3/(1+x^4)^2\,\dd x$。
\end{exercise}
\begin{exercise}{527}{常规·四次分母典型模型}
求 $\int(x^2+1)/(x^4+3x^2+2)\,\dd x$。提示：先分解分母并约分。
\end{exercise}
\begin{exercise}{528}{基础·原函数奇偶性实例}
求 $f(x)=x^3$ 满足 $F(0)=0$ 的原函数，并判断 $F$ 的奇偶性。
\end{exercise}
\begin{exercise}{529}{基础·连续函数原函数}
设 $f(x)=x^2+1$。求满足 $F(0)=0$ 的原函数，并说明 $f$ 在 $[-1,1]$ 上有最大、最小值的依据。
\end{exercise}
\begin{exercise}{530}{常规·分段原函数参数}
设 $f(x)=2(x-1)$（$x<1$）、$f(x)=\ln x$（$x\ge1$）。求 $a,b$，使 $F(x)=(x-1)^2$（$x<1$）、$F(x)=x(\ln x+a)+b$（$x\ge1$）为 $f$ 在实轴上的原函数。
\end{exercise}
\begin{exercise}{531}{常规·反正切指数分部}
求 $\int e^x\arctan(e^x)\,\dd x$。提示：先令 $u=e^x$，再分部积分。
\end{exercise}
\begin{exercise}{532}{基础·反正切整体换元}
求 $\int\arctan x/(1+x^2)\,\dd x$。
\end{exercise}
\begin{exercise}{533}{常规·指数有理对数导数}
求 $\int(1+x)/[x(1+xe^x)]\,\dd x$（$x>0$）。提示：将被积函数写成 $1+1/x-(xe^x)'/(1+xe^x)$。
\end{exercise}
\begin{exercise}{534}{常规·对数商逆向求导}
求 $\int(1-\ln x)/(x-\ln x)^2\,\dd x$（$x>0$）。可使用 $x-\ln x\ge1$，并计算 $[x/(x-\ln x)]'$。
\end{exercise}
\begin{exercise}{535}{常规·根式反正切换元}
求 $\int e^x\arctan\sqrt{e^x-1}/\sqrt{e^x-1}\,\dd x$（$x>0$）。提示：令 $t=\sqrt{e^x-1}$，积分化为 $2\int\arctan t\,\dd t$。
\end{exercise}
\begin{exercise}{536}{常规·商函数导数结构}
求 $\int[x+\sin x\cos x]/(x\sin x+\cos x)^2\,\dd x$，在分母非零的开区间上考虑。提示：计算 $[\cos x/(x\sin x+\cos x)]'$。
\end{exercise}
\begin{exercise}{537}{基础·三角有理换元}
求 $\int\cos x/[\sin x(1+\sin x)]\,\dd x$，在 $\sin x\ne0,-1$ 的开区间上考虑。提示：令 $u=\sin x$。
\end{exercise}
\begin{exercise}{538}{常规·万能代换}
在 $(-\pi/2,\pi/2)$ 上求 $\int\dd x/(1+\sin x+\cos x)$。提示：令 $t=\tan(x/2)$。
\end{exercise}
\begin{exercise}{539}{常规·万能代换与反正切}
在 $(-\pi,\pi)$ 上求 $\int\dd x/(2+\cos x)$。提示：令 $t=\tan(x/2)$。
\end{exercise}
\begin{exercise}{540}{基础·分母及其导数拆分}
求 $\int(3\cos x-\sin x)/(\cos x+\sin x)\,\dd x$，在分母非零的开区间上考虑。
\end{exercise}
\begin{exercise}{541}{基础·对称拆分三角积分}
求 $\int\sin x/(\sin x+\cos x)\,\dd x$，在分母非零的开区间上考虑。
\end{exercise}
\begin{exercise}{542}{常规·根式三角代换}
求 $\int\dd x/[2x+\sqrt{1-x^2}]$，在 $0<x<1$ 上考虑。提示：令 $x=\sin t$（$0<t<\pi/2$），把所得分子写成分母及其导数的线性组合。
\end{exercise}
\begin{exercise}{543}{常规·拆项与重二次分母}
求 $\int(x^3+1)/(1+x^2)^2\,\dd x$。提示：分开处理 $x^3/(1+x^2)^2$ 与 $1/(1+x^2)^2$。
\end{exercise}
\begin{exercise}{544}{基础·反正切指数换元}
求 $\int e^{\arctan x}/(1+x^2)\,\dd x$。
\end{exercise}
\begin{exercise}{545}{常规·指数有理换元}
求 $\int\dd x/(1+e^{x/2})$。提示：令 $t=e^{x/2}$，注意 $\dd x=2\dd t/t$。
\end{exercise}
\begin{exercise}{546}{常规·三角因子有理换元}
在 $0<x<\pi$ 上求 $\int\sin x/(2+\cos x)\,\dd x$。
\end{exercise}
\begin{exercise}{547}{基础·反正弦整体换元}
求 $\int\arcsin x/\sqrt{1-x^2}\,\dd x$（$|x|<1$）。
\end{exercise}
\begin{exercise}{548}{常规·已知函数关系的积分}
设 $y=y(x)>0$ 满足 $y^2=x$（$x>0$）。求 $\int\dd x/y(x)$，并用 $x=y^2$ 换元核验。
\end{exercise}
\begin{exercise}{549}{常规·函数关系与换元}
已知 $f(\sin^2t)=t/\sin t$（$0<t<\pi/2$）。求 $\int f(x)/\sqrt{1-x}\,\dd x$（$0<x<1$）。提示：令 $x=\sin^2t$，积分化为 $\int2t\,\dd t$。
\end{exercise}
\begin{exercise}{550}{常规·已知原函数与分部积分}
已知 $\int f(x)\,\dd x=e^{x^2}+C$。求 $f(x)$ 及 $\int xf'(x)\,\dd x$。提示：分部积分可写为 $xf(x)-\int f(x)\,\dd x$。
\end{exercise}
\begin{exercise}{551}{基础·加权原函数恢复}
已知 $\int xf(x)\,\dd x=x^3+C$（$x>0$）。求 $f(x)$ 和 $\int f(x)\,\dd x$。
\end{exercise}
\begin{exercise}{552}{常规·复合导函数恢复}
已知 $f'(e^t)=te^{-t}$ 对所有实数 $t$ 成立且 $f(1)=0$。求 $f(x)$（$x>0$）。
\end{exercise}
\begin{exercise}{553}{常规·抽象函数平方导数}
设 $f\in C^1(I)$。求 $\int f(x)f'(x)\,\dd x$ 及 $\int f'(x)/(1+f(x)^2)\,\dd x$。
\end{exercise}
\begin{exercise}{554}{常规·线性三角组合平方}
求 $\int\dd x/(\sin x+\cos x)^2$，在分母非零的开区间上考虑。提示：先写成 $2\sin^2(x+\pi/4)$。
\end{exercise}
\begin{exercise}{555}{基础·部分分式变式}
求 $\int\dd x/[(x-1)(x-2)]$，在不含 $1,2$ 的开区间上考虑。
\end{exercise}
\begin{exercise}{556}{常规·有理积分小综合}
求 $\int\dd x/[x(x^2+1)^2]$，在两个半轴上分别考虑。提示：先拆为 $1/x-x/(x^2+1)-x/(x^2+1)^2$。
\end{exercise}
